\section{结\quad 论}

本文以溶胶-凝胶与静电纺丝两条路线制备的 TiO$_2$ 为对象，用 XRD、拉曼光谱、SEM/EDS 与紫外-可见漫反射对同一温度序列（200—900 $^\circ$C）的产物进行了系统表征，主要结论如下。

(1) 相变过程具有明确的温度阈值与 S 形特征。200 $^\circ$C 样品近于非晶；300—600 $^\circ$C 为纯锐钛矿；700 $^\circ$C 金红石质量分数突增至 17.3\%，800 $^\circ$C 达 84.7\%，900 $^\circ$C 完全转变为金红石。锐钛矿 (101) 峰的半高宽由 1.19$^\circ$ 降至 0.18$^\circ$，Scherrer 晶粒尺寸由 6.8 nm 增至 44.8 nm，且 500 $^\circ$C 之后长大速率明显加快。

(2) 拉曼光谱独立验证了物相转变，并揭示相变在微区上的不均匀性。锐钛矿的 144、399、516、639 cm$^{-1}$ 峰随温度升高而减弱，金红石的 448 与 611 cm$^{-1}$ 峰逐渐出现并增强；700 $^\circ$C 的两个平行样品分别呈纯锐钛矿与纯金红石，说明此时相变尚未在整个样品内均匀完成。

(3) 形貌演化与晶粒生长同步。SEM 统计得到 300、500、700 $^\circ$C 颗粒平均尺寸为 31$\pm$11、39$\pm$12、85$\pm$28 nm，900 $^\circ$C 粗化至约 0.9 $\mu$m 并出现烧结颈；颗粒尺寸约为 XRD 晶粒尺寸的 1.5—2.5 倍，表明颗粒多为多晶聚集体。静电纺丝产物为直径中位数约 0.48 $\mu$m 的连续纤维，伴少量珠状缺陷。

(4) 光学带隙随相变与晶粒长大单调下降。由 Tauc 关系外推得 400—600 $^\circ$C 的表观带隙为 3.12—3.13 eV，700 $^\circ$C 为 3.05 eV，800 $^\circ$C 为 2.80 eV，与金红石含量增加、量子尺寸效应减弱相一致。EDS 仅检出 Ti、O、C 与基底 Al，未发现其他杂质。

(5) 制备路线主要决定产物的宏观组装形态（颗粒堆积体或连续纤维膜），而热处理温度决定相组成、晶粒尺寸与带隙。因此若要调控 TiO$_2$ 的结构与光学性质，应以热处理制度为首要变量，再根据应用需求选择成型方式。

通过本课程的学习与本次论文的数据整理，笔者体会到：单一表征只能给出局部信息，只有把衍射（长程平均）、拉曼（短程有序）、显微（形貌与尺寸）与光谱（电子结构）的结果放在同一坐标系下交叉验证，才能形成对材料“结构—性能”关系的完整认识；同时，样品制备与装样的一致性对定量结果的影响往往被低估。今后希望能补充 Rietveld 精修、BET 比表面积测试与光催化降解实验，把结构表征与性能评价真正打通。
